\section{Model}
\label{sec:model}

There are two symmetric regions, indexed by $i\in\{1,2\}$, each governed by a local government. Each government controls a fixed unit of industrial-policy capacity. It allocates a share $x_i\in[0,1]$ to a locally high-return activity, denoted $U$, and the remaining share $1-x_i$ to a complementary activity, denoted $D$. The fixed-capacity restriction is substantive: increasing support for one activity necessarily reduces support for the other. The paper therefore studies the composition of policy rather than its total level.

\subsection{Direct local returns}

Let $b_D$ denote the direct real surplus generated by a unit allocated to $D$, and let
\begin{equation}
\Delta\equiv b_U-b_D>0
\label{eq:delta}
\end{equation}
be the direct local surplus premium from allocating capacity to $U$. Thus the direct component of region $i$'s payoff is
\begin{equation}
b_D+\Delta x_i.
\label{eq:direct-return}
\end{equation}
The parameter $\Delta$ is an opportunity-cost-adjusted real surplus premium, such as a productivity, real-value-added, or real-wage advantage. It is not interpreted as a pure fiscal transfer or political rent.

\subsection{Cross-regional complementarity}

The two activities are vertically complementary across regions. A matched pair consisting of one unit of $U$ and one unit of $D$ produces
\begin{equation}
Y=A u^{\alpha}d^{1-\alpha},
\qquad A>0,\quad 0<\alpha<1.
\label{eq:production}
\end{equation}
At the unit match $(u,d)=(1,1)$, total output is $A$. Under competitive constant-returns incidence, the $U$ side receives $\alpha A$ and the $D$ side receives $(1-\alpha)A$. Hence $\alpha$ determines how the complementary surplus is distributed across the two jurisdictions; it is not an arbitrary transfer share imposed after production.

Projects are matched across regions. Given portfolio choices $(x_i,x_j)$, the mass of matches with $U$ in region $i$ and $D$ in region $j$ is $x_i(1-x_j)$, while the reverse-direction mass is $(1-x_i)x_j$. Region $i$ therefore receives $\alpha A$ on the former matches and $(1-\alpha)A$ on the latter.

Combining direct and cross-regional returns, the frozen local payoff is
\begin{equation}
W_i(x_i,x_j)
=b_D+\Delta x_i
+\alpha A x_i(1-x_j)
+(1-\alpha)A(1-x_i)x_j,
\qquad j\neq i.
\label{eq:local-payoff}
\end{equation}
Governments choose simultaneously under complete information and maximize their own regional payoff $W_i$. The equilibrium concept is pure-strategy Nash equilibrium on $[0,1]^2$. Because strategies are continuous portfolio shares, an interior equilibrium below is a pure-strategy equilibrium, not a mixed equilibrium.

The baseline deliberately abstracts from endogenous policy budgets, firm pricing, entry, migration, and tax financing. Those margins are not needed to isolate the composition-specific decentralization problem studied here.
